\chapter{Hierarchical latents}
\label{ch:hier}

One latent layer can be too weak a prior. A two-layer model draws a coarse code and then a fine code:
\[
p_\theta(x,z_1,z_2)=p_\theta(z_2)\,p_\theta(z_1\given z_2)\,p_\theta(x\given z_1),
\]
with the conditional independence $x\perp z_2\given z_1$. The variational family used in the homework respects the same order,
\[
q_\phi(z_1,z_2\given x)=q_\phi(z_2\given x)\,q_\phi(z_1\given z_2,x).
\]

\begin{center}
\begin{tikzpicture}[
  lat/.style={circle, draw=burgundy, thick, minimum size=9mm, fill=burgundysoft},
  obs/.style={circle, draw=navy, thick, minimum size=9mm, fill=navysoft},
  arr/.style={-{Stealth[length=2.2mm]}, thick, draw=slate}
]
\node[lat] (z2) {$z_2$};
\node[lat, below=11mm of z2] (z1) {$z_1$};
\node[obs, below=11mm of z1] (x) {$x$};
\draw[arr] (z2) -- (z1);
\draw[arr] (z1) -- (x);
\end{tikzpicture}
\end{center}

\section{The hierarchical ELBO}

Apply~\eqref{eq:elbokl} to the pair $(z_1,z_2)$. The gap is
\[
\log p_\theta(x)-\mathcal L_H(x)
=\KL\bigl(q(z_1,z_2\given x)\Vert p(z_1,z_2\given x)\bigr).
\]
Expand the joint and the variational factorizations. The cross terms separate into three pieces:
\begin{align*}
\mathcal L_H(x)
&=\E_{q(z_1,z_2\given x)}[\log p(x\given z_1)]
-\KL\bigl(q(z_2\given x)\Vert p(z_2)\bigr)\\
&\quad-\E_{q(z_2\given x)}\Bigl[\KL\bigl(q(z_1\given z_2,x)\Vert p(z_1\given z_2)\bigr)\Bigr].
\end{align*}
The inner KL is averaged under $q(z_2\given x)$ because $z_2$ is still random when that conditional KL is evaluated. You draw the coarse code from the approximate posterior, then score the fine code.

\section{Comparison with a one-layer bound}

Let $q(z_1\given x)$ be the marginal of the variational joint, and define the ordinary ELBO that uses only $z_1$,
\[
\mathcal L_M(x)
=\E_{q(z_1\given x)}[\log p(x\given z_1)]
-\KL\bigl(q(z_1\given x)\Vert p(z_1)\bigr).
\]
A second KL decomposition, this time of $q(z_2\given z_1,x)$ against $p(z_2\given z_1)$, yields
\[
\mathcal L_H(x)
=\mathcal L_M(x)
-\E_{q(z_1\given x)}\Bigl[\KL\bigl(q(z_2\given z_1,x)\Vert p(z_2\given z_1)\bigr)\Bigr].
\]
The subtracted term is nonnegative, so $\mathcal L_H\le\mathcal L_M$. Equality holds if and only if $q(z_2\given z_1,x)=p(z_2\given z_1)$ almost everywhere. For a fixed marginal $q(z_1\given x)$, that conditional is exactly the one that maximizes $\mathcal L_H$.

The marginal prior $p(z_1)=\int p(z_1\given z_2)p(z_2)\,dz_2$ may have no closed form. For an arbitrary positive auxiliary $r(z_2\given z_1)$,
\begin{align*}
\log p(z_1)
&=\E_{r}[\log p(z_1\given z_2)]-\KL(r(z_2\given z_1)\Vert p(z_2))\\
&\quad+\KL(r(z_2\given z_1)\Vert p(z_2\given z_1)).
\end{align*}
Drop the last KL, which is nonnegative, and you have a lower bound. The bound is an equality when $r=p(z_2\given z_1)$. Therefore
\[
\log p(z_1)
=\sup_r\Bigl\{\E_r[\log p(z_1\given z_2)]-\KL(r\Vert p(z_2))\Bigr\},
\]
and the supremum is attained at the true conditional. This is the same ``gap equals a KL'' pattern as~\eqref{eq:elbokl}.

\begin{practbox}
In the decomposition of $\mathcal L_H$, why is the fine-code KL inside an expectation over $q(z_2\given x)$, rather than evaluated at a single $z_2$?
\end{practbox}
\begin{solbox}
The KL $\KL(q(z_1\given z_2,x)\Vert p(z_1\given z_2))$ is a function of $z_2$. The variational joint draws $z_2$ from $q(z_2\given x)$ before drawing $z_1$. Taking the expectation is what the joint KL produces when you integrate that coordinate. Freezing one $z_2$ would describe a different, degenerate inference distribution.
\end{solbox}
